\documentclass[12pt,a4paper]{article}
\usepackage[UTF8]{ctex}
\usepackage{amsmath,amssymb,amsthm}
\usepackage{geometry}
\usepackage{graphicx}

\geometry{margin=2.5cm}

\title{二阶系统阻尼比选择：临界阻尼是否最好？}
\author{第11章补充说明}
\date{}

\begin{document}

\maketitle

\section{问题}

对于二阶误差动力学系统：
\begin{equation}
\ddot{\theta}_e(t) + 2\zeta\omega_n \dot{\theta}_e(t) + \omega_n^2 \theta_e(t) = 0,
\end{equation}
其中$\zeta$是阻尼比，$\omega_n$是自然频率。

\textbf{问题：临界阻尼（$\zeta = 1$）是否是最好的选择？}

\section{三种阻尼情况回顾}

\subsection{过阻尼：$\zeta > 1$}

特征根为两个不同的负实数：
\begin{equation}
s_{1,2} = -\zeta\omega_n \pm \omega_n\sqrt{\zeta^2 - 1}
\end{equation}

解的形式：
\begin{equation}
\theta_e(t) = c_1 e^{s_1 t} + c_2 e^{s_2 t}
\end{equation}

\textbf{特点：}
\begin{itemize}
\item 无超调，响应单调递减
\item 响应速度慢（由较慢的时间常数$\tau_1 = -1/s_1$决定）
\item 稳定时间较长
\end{itemize}

\subsection{临界阻尼：$\zeta = 1$}

特征根为两个相等的负实数：
\begin{equation}
s_1 = s_2 = -\omega_n
\end{equation}

解的形式：
\begin{equation}
\theta_e(t) = (c_1 + c_2 t)e^{-\omega_n t}
\end{equation}

对于单位误差响应（$\theta_e(0) = 1$，$\dot{\theta}_e(0) = 0$）：
\begin{equation}
\theta_e(t) = (1 + \omega_n t)e^{-\omega_n t}
\end{equation}

\textbf{特点：}
\begin{itemize}
\item 无超调，响应单调递减
\item 在无超调的情况下，收敛速度最快
\item 稳定时间：约$4.7/\omega_n$（2\%稳定时间）
\end{itemize}

\subsection{欠阻尼：$\zeta < 1$}

特征根为复共轭：
\begin{equation}
s_{1,2} = -\zeta\omega_n \pm j\omega_n\sqrt{1 - \zeta^2} = -\zeta\omega_n \pm j\omega_d
\end{equation}
其中$\omega_d = \omega_n\sqrt{1 - \zeta^2}$是阻尼自然频率。

解的形式：
\begin{equation}
\theta_e(t) = e^{-\zeta\omega_n t}(c_1 \cos(\omega_d t) + c_2 \sin(\omega_d t))
\end{equation}

对于单位误差响应：
\begin{equation}
\theta_e(t) = e^{-\zeta\omega_n t}\left(\cos(\omega_d t) + \frac{\zeta\omega_n}{\omega_d}\sin(\omega_d t)\right)
\end{equation}

\textbf{特点：}
\begin{itemize}
\item 有超调（除非$\zeta$很大）
\item 响应速度快（可能比临界阻尼更快）
\item 有振荡，但会衰减
\end{itemize}

\section{性能指标比较}

\subsection{超调量}

超调量（Overshoot）定义为误差响应首次达到最小值时的相对值：
\begin{equation}
\text{超调} = \frac{\theta_{e,\min} - e_{ss}}{\theta_e(0) - e_{ss}} \times 100\%
\end{equation}

对于欠阻尼系统（$\zeta < 1$），超调量为：
\begin{equation}
\text{超调} = e^{-\pi\zeta/\sqrt{1-\zeta^2}} \times 100\%
\end{equation}

\begin{itemize}
\item $\zeta = 1$：超调 = 0\%
\item $\zeta = 0.7$：超调 $\approx$ 4.6\%
\item $\zeta = 0.5$：超调 $\approx$ 16.3\%
\item $\zeta = 0.3$：超调 $\approx$ 37.2\%
\end{itemize}

\subsection{稳定时间}

2\%稳定时间$T_s$是误差首次进入并保持在$[0.02\theta_e(0), -0.02\theta_e(0)]$范围内的时间。

对于欠阻尼系统，近似公式为：
\begin{equation}
T_s \approx \frac{4}{\zeta\omega_n}
\end{equation}

对于临界阻尼（$\zeta = 1$）：
\begin{equation}
T_s \approx \frac{4.7}{\omega_n}
\end{equation}

\textbf{比较：}
\begin{itemize}
\item $\zeta = 1$：$T_s \approx 4.7/\omega_n$
\item $\zeta = 0.7$：$T_s \approx 5.7/\omega_n$（稍慢）
\item $\zeta = 0.5$：$T_s \approx 8/\omega_n$（更慢）
\end{itemize}

\subsection{上升时间}

上升时间$T_r$是误差从10\%衰减到90\%所需的时间。

对于欠阻尼系统：
\begin{equation}
T_r \approx \frac{1.8}{\omega_n} \quad \text{（当$\zeta \approx 0.7$时）}
\end{equation}

对于临界阻尼：
\begin{equation}
T_r \approx \frac{2.2}{\omega_n}
\end{equation}

\textbf{结论：}欠阻尼系统的上升时间通常更短。

\subsection{峰值时间}

峰值时间$T_p$是误差达到最小值（最大超调）的时间。

对于欠阻尼系统：
\begin{equation}
T_p = \frac{\pi}{\omega_d} = \frac{\pi}{\omega_n\sqrt{1-\zeta^2}}
\end{equation}

临界阻尼系统没有峰值（无超调）。

\section{为什么临界阻尼不一定最好？}

\subsection{关键权衡：超调 vs 响应速度}

这是一个\textbf{权衡取舍}的问题：

\begin{itemize}
\item \textbf{临界阻尼（$\zeta = 1$）：}无超调，但在无超调前提下收敛最快
\item \textbf{欠阻尼（$\zeta = 0.7$）：}有超调（约4.6\%），但响应更快
\end{itemize}

\textbf{重要说明：}如果应用\textbf{绝对不能有超调}，那么临界阻尼确实是最好的选择。但在许多应用中，可以接受少量超调来换取更快的响应速度。

\subsection{理论分析}

虽然临界阻尼在\textbf{无超调}的前提下提供了最快的收敛速度，但在实际应用中，我们通常可以接受\textbf{少量超调}来换取\textbf{更快的响应速度}。

\subsection{最优阻尼比：$\zeta = 0.7$（在可接受超调的前提下）}

\textbf{注意：}$\zeta = 0.7$确实有超调（约4.6\%）！

在控制工程实践中，$\zeta = 0.7$（约$1/\sqrt{2}$）通常被认为是最优选择，\textbf{前提是应用可以接受少量超调}，原因如下：

\begin{enumerate}
\item \textbf{快速响应：}上升时间短（比临界阻尼快约22\%），初始响应快
\item \textbf{适度超调：}超调量约4.6\%，在许多应用中可接受
\item \textbf{良好的稳定时间：}虽然比临界阻尼稍慢，但差异不大（约21\%）
\item \textbf{鲁棒性：}对参数变化不敏感
\item \textbf{综合性能最优：}在多个性能指标之间取得良好平衡（如ISE、ITAE等积分指标）
\end{enumerate}

\textbf{但如果应用不能容忍超调，临界阻尼（$\zeta = 1$）就是更好的选择。}

\subsection{性能指标优化}

考虑积分性能指标：

\textbf{ISE（积分平方误差）：}
\begin{equation}
\text{ISE} = \int_0^{\infty} \theta_e^2(t) dt
\end{equation}

\textbf{ITAE（积分时间加权绝对误差）：}
\begin{equation}
\text{ITAE} = \int_0^{\infty} t|\theta_e(t)| dt
\end{equation}

对于这些指标，最优阻尼比通常在$\zeta = 0.6$到$\zeta = 0.8$之间，而不是$\zeta = 1$。

\section{不同应用场景的选择}

\subsection{需要无超调的应用}

如果系统\textbf{绝对不能有超调}（例如，避免碰撞、保护设备），则：
\begin{itemize}
\item 选择临界阻尼（$\zeta = 1$）或轻微过阻尼（$\zeta = 1.1 \sim 1.2$）
\item 牺牲响应速度换取安全性
\end{itemize}

\subsection{追求快速响应的应用}

如果系统需要\textbf{快速响应}，可以接受少量超调：
\begin{itemize}
\item 选择$\zeta = 0.7$（最优折衷）
\item 或$\zeta = 0.6 \sim 0.8$（根据具体需求调整）
\end{itemize}

\subsection{需要平滑响应的应用}

如果系统需要\textbf{平滑、无振荡}的响应：
\begin{itemize}
\item 选择过阻尼（$\zeta > 1$，如$\zeta = 1.5 \sim 2.0$）
\item 响应会很慢，但非常平滑
\end{itemize}

\subsection{机器人控制应用}

在机器人控制中：
\begin{itemize}
\item \textbf{位置控制：}通常选择$\zeta = 0.7 \sim 1.0$，平衡速度和稳定性
\item \textbf{力控制：}可能需要$\zeta = 1.0 \sim 1.2$，避免力超调
\item \textbf{轨迹跟踪：}$\zeta = 0.7$通常是最佳选择
\end{itemize}

\section{数值示例}

考虑$\omega_n = 1$ rad/s的二阶系统，比较不同阻尼比的响应：

\subsection{$\zeta = 1.5$（过阻尼）}

\begin{itemize}
\item 稳定时间：$\approx 6.5/\omega_n = 6.5$ s
\item 超调：0\%
\item 响应：非常慢，但平滑
\end{itemize}

\subsection{$\zeta = 1.0$（临界阻尼）}

\begin{itemize}
\item 稳定时间：$\approx 4.7/\omega_n = 4.7$ s
\item 超调：0\%
\item 响应：在无超调情况下最快
\end{itemize}

\subsection{$\zeta = 0.7$（最优欠阻尼）}

\begin{itemize}
\item 稳定时间：$\approx 5.7/\omega_n = 5.7$ s
\item 超调：4.6\%
\item 上升时间：$\approx 1.8/\omega_n = 1.8$ s（比临界阻尼快22\%）
\item 响应：快速，轻微超调
\end{itemize}

\subsection{$\zeta = 0.5$（欠阻尼）}

\begin{itemize}
\item 稳定时间：$\approx 8.0/\omega_n = 8.0$ s
\item 超调：16.3\%
\item 响应：有振荡，稳定时间较长
\end{itemize}

\section{结论}

\textbf{临界阻尼（$\zeta = 1$）vs 欠阻尼（$\zeta = 0.7$）：这是一个权衡取舍的问题。}

\begin{enumerate}
\item \textbf{临界阻尼（$\zeta = 1$）的优势：}
   \begin{itemize}
   \item \textbf{无超调}（这是最重要的优势！）
   \item 在无超调前提下收敛最快
   \item 响应单调递减，平滑
   \end{itemize}

\item \textbf{临界阻尼（$\zeta = 1$）的劣势：}
   \begin{itemize}
   \item 上升时间不是最快的（比$\zeta = 0.7$慢约22\%）
   \item 对于某些性能指标（如ISE、ITAE）不是最优
   \end{itemize}

\item \textbf{欠阻尼（$\zeta = 0.7$）的优势：}
   \begin{itemize}
   \item 上升时间更快（比临界阻尼快约22\%）
   \item 对于某些积分性能指标更优
   \end{itemize}

\item \textbf{欠阻尼（$\zeta = 0.7$）的劣势：}
   \begin{itemize}
   \item \textbf{有超调}（约4.6\%）
   \item 稳定时间稍慢（比临界阻尼慢约21\%）
   \end{itemize}

\item \textbf{推荐选择：}
   \begin{itemize}
   \item \textbf{不能有超调的应用：} $\zeta = 1.0$（临界阻尼）\textbf{是最好的选择}
   \item \textbf{可以接受少量超调的一般应用：} $\zeta = 0.7$（综合性能更好）
   \item \textbf{需要极快响应：} $\zeta = 0.6 \sim 0.7$（可接受更大超调）
   \item \textbf{需要平滑响应：} $\zeta = 1.2 \sim 2.0$（过阻尼，无超调但更慢）
   \end{itemize}
\end{enumerate}

\textbf{关键点：}
\begin{itemize}
\item 如果应用\textbf{不能容忍任何超调}，临界阻尼（$\zeta = 1$）\textbf{就是最好的选择}。
\item 如果应用\textbf{可以接受少量超调}（如4.6\%），$\zeta = 0.7$通常能提供更好的综合性能（更快的响应速度）。
\item 最优阻尼比取决于具体的应用需求、性能指标和约束条件。
\end{itemize}

\section{设计建议}

在实际控制系统设计中：

\begin{enumerate}
\item \textbf{首先判断：能否接受超调？}
   \begin{itemize}
   \item \textbf{不能接受超调：}直接选择$\zeta = 1.0$（临界阻尼）或$\zeta > 1.0$（过阻尼）
   \item \textbf{可以接受超调：}继续下一步
   \end{itemize}

\item \textbf{如果可以接受超调，从$\zeta = 0.7$开始：}这是大多数应用的良好起点
\item \textbf{根据需求调整：}
   \begin{itemize}
   \item 如果超调不可接受，增加$\zeta$到$1.0$（临界阻尼）或更高（过阻尼）
   \item 如果需要更快响应，可以降低$\zeta$到$0.6$（但要注意超调增加到约9.5\%）
   \item 如果需要更平滑响应，增加$\zeta$到$1.2$以上（过阻尼，无超调但更慢）
   \end{itemize}
\item \textbf{考虑鲁棒性：}$\zeta = 0.7$对参数变化不敏感，鲁棒性好；$\zeta = 1.0$也有良好的鲁棒性
\item \textbf{仿真验证：}在实际应用前，通过仿真验证不同$\zeta$值的性能，特别是超调量
\end{enumerate}

\textbf{决策流程：}
\begin{enumerate}
\item 应用是否\textbf{绝对不能有超调}？
   \begin{itemize}
   \item 是 $\rightarrow$ 选择$\zeta = 1.0$（临界阻尼）或$\zeta > 1.0$（过阻尼）
   \item 否 $\rightarrow$ 继续
   \end{itemize}
\item 可以接受多少超调？
   \begin{itemize}
   \item 可接受4-5\% $\rightarrow$ 选择$\zeta = 0.7$
   \item 可接受更多 $\rightarrow$ 可以选择$\zeta = 0.6$或更小
   \item 希望更保守 $\rightarrow$ 选择$\zeta = 0.8 \sim 0.9$
   \end{itemize}
\item 通过仿真验证最终选择
\end{enumerate}

\end{document}

